%!TEX root = ../kr.tex
\chapter{Маркетинг чи обман? Межа споживчої пильності}
\label{apdx:marketing_logic}

% Компактний шрифт для економії місця
\small
\setlength{\parskip}{3pt}

\textit{Маркетинг грає з вашими бажаннями, а шахрайство — з вашою безпекою.}


\section*{Частина 1. Де закінчується законний маркетинг?}
\vspace{-0.1cm}

\textbf{Маркетинг (Законно, але маніпулятивно):}
\begin{itemize}[leftmargin=2em, itemsep=0pt, parsep=3pt]
    \item \textbf{Ціна за одиницю, вагу.} Ви бачите низьку ціну, але дрібним шрифтом вказано, що це за 100 грамів, 1 штуку або за умови купівлі трьох одиниць. Це гра на вашій неуважності.
    \item \textbf{Шринкфляція.} Зменшення об'єму пакування (наприклад, 850 мл замість 1 л.) при збереженні ціни та дизайну. Ви отримуєте менше товару, але закон не порушено.
    \item \textbf{Фуд-стилістика.} Ретушовані фото у меню чи на коробці. Якщо склад і вага відповідають заявленим, то ідеальний вигляд на фото — це лише рекламний трюк, а не обман.
\end{itemize}

\textbf{Шахрайство (Злочин):}
\begin{itemize}[leftmargin=2em, itemsep=0pt, parsep=3pt]
    \item \textbf{Свідома дезінформація.} У характеристиках вказано «натуральна шкіра», а приходить дешевий замінник. Або вказано «новий товар», а приходить вживаний.
    \item \textbf{Приховування дефектів.} Продавець навмисно робить фото так, щоб приховати тріщини, сколи або несправність електроніки.
    \item \textbf{Фантомний товар.} Використання професійних фото з інтернету для продажу товару, якого у продавця взагалі немає в наявності.
\end{itemize}

 

\section*{Частина 2. Контрольні питання для перевірки «обкладинки»}
\vspace{-0.1cm}



\subsection*{1. Психологія «ідеального фото»}
Маркетолог хоче, щоб ви захотіли товар через емоцію. Шахрай хоче, щоб ви втратили пильність через нереально низьку ціну.

\textbf{Подумайте!} Чи відповідає ціна ринковій? Якщо iPhone або брендова сумка коштують у три рази дешевше, ніж усюди — жоден «розпродаж» не виправдовує таку ціну. Ви купуєте не товар, а ризик.

\vspace{0.25cm}

\subsection*{2. Пастка «дрібного шрифту»}
У магазині ми часто звинувачуємо касира: «На ціннику було інше!». Але часто ми просто не дочитали умови акції. 

\textbf{Подумайте!} Чи не переносите ви цю неуважність у цифровий простір? Шахраї часто використовують схожу тактику, ховаючи умови «автоматичного списання коштів» або «підписки» у довгих текстах, які ніхто не читає.

 

\section*{Висновок. Як не дати себе обманути?}
\vspace{-0.1cm}

Межа між «хитрим продавцем» та «злочинцем» проходить там, де ви втрачаєте контроль над своїми грошима. 

 

\textbf{Три золоті правила споживача:}
\begin{enumerate}[itemsep=3pt, parsep=0pt]
    \item \textbf{Рахуйте вартість за одиницю.} Завжди переводьте ціну у стандартні одиниці (кілограми, літри), щоб зрозуміти реальну вигоду. 
    \item \textbf{Запитуйте «живі» фото.} Якщо купуєте онлайн на платформах оголошень — просіть фото товару з папірцем, на якому написано сьогоднішню дату. Шахрай не зможе цього зробити.
    \item \textbf{Читайте умови повернення.} Маркетолог завжди дасть можливість повернути товар (згідно із законом), шахрай — зробить усе, щоб ви не змогли його знайти після оплати.
\end{enumerate}


\newpage
